\subsection{The benchmark marks a correct answer wrong, which bounds the cost figures above}
\label{sec:benchmark}

No score in Sections~\ref{sec:controlled} to~\ref{sec:hard} is sounder than its ground truth, and on
eICU it is wrong in both directions. It contains a body temperature of 99.8 degrees Celsius, a
Fahrenheit value labelled with a Celsius unit: real corruption that nothing injected. A truthful
bound flags it, and the scoring counts that correct detection as a false positive, since the ground
truth assumes every uninjected record clean. It is the only \texttt{known\_corrupt\_values} record
(Section~\ref{sec:fair}) inside the 200-resource evaluation slice.

Relabelling that allowlist as genuine defects lifts the twelve-cell mean plausibility F1 from
0.9851 to 0.9908, with 7 of 12 cells perfect. That relabelling has \textbf{not} been applied,
because redefining ground truth after seeing the scores manufactures a result.

The second direction has no allowlist and no counterfactual, and nothing measures what the four such
cases cost. Where the legacy injector writes 9999 into an analyte whose ceiling is legitimately
above 9999, the injected value is credible and a truthful bound passes it, scored as a miss.

F1 therefore carries an unknown ground-truth error in both directions. The bound falls on the cost
half alone: the privacy half is counted in released fields and never scored against injected ground
truth. Appendix E documents how the error surfaced, the rest of the allowlist, and the four cases.
